\begin{frame}{활성함수가 왜 필요한가}
  \begin{itemize}
    \item 활성함수가 없으면: $a_2 = W_2(W_1 x) = (W_2 W_1) x$ —
      \textbf{층을 아무리 쌓아도 한 번의 선형 변환}과 같음
    \item 9.1에서 배운 것처럼 선형 모델로는 XOR이 \textbf{원칙적으로
      불가} — 활성함수 없이는 100층을 쌓아도 영원히 못 풀음
    \item 활성함수는 각 층 사이에 ``곡선을 접어 넣는'' \kb{비선형성}을
      제공, 층이 쌓일수록 표현 가능 공간이 \textbf{지수적으로} 넓어짐
  \end{itemize}
  \vskip0.4em
  {\footnotesize 1958년 퍼셉트론의 스텝함수($z \ge 0$이면 1, 아니면 0)가
  역사적으로 가장 첫 번째 ``활성함수'' — 이 맥락에서 보면 자연스러움.
  본 절의 활성함수·초기화와 9.3의 정규화는 전부 ``깊은 망에서도
  그래디언트가 살아서 전달되게 하려면''이라는 하나의 문제의
  서로 다른 각도의 답이다.}
\end{frame}
